\documentclass[11pt,a4paper]{article}

\usepackage[margin=2.15cm]{geometry}
\usepackage{amsmath,amssymb,amsfonts}
\usepackage{bm}
\usepackage{booktabs}
\usepackage{array}
\usepackage{hyperref}
\usepackage{xcolor}
\usepackage{tikz}
\usepackage{pgfplots}
\usepackage{subcaption}
\usepackage{float}

\pgfplotsset{compat=1.18}
\usetikzlibrary{arrows.meta,positioning,calc}

\hypersetup{
    colorlinks=true,
    linkcolor=blue!60!black,
    citecolor=blue!60!black,
    urlcolor=blue!60!black
}

\title{Mathematical Derivation of the TwoTimescaleResilience Framework:\\
Species-Aware Pipeline with Exposure and Mode of Action Mapping}
\author{Model Documentation}
\date{\today}

\begin{document}
\maketitle

\begin{abstract}
This document provides a full mathematical derivation of the \texttt{TwoTimescaleResilience.jl} spatial framework. It explicitly details how raw environmental rasters are translated into species-specific vulnerabilities. By introducing a target group index $(q)$, we trace the pipeline from raw ambient stress, through species-specific exposure filters, intermediate Modes of Action (MoA), and finally to Dynamic Energy Budget (DEB) physiological axes. The derivation concludes with the core integration theorem, proving that amplification of a perturbation response is inversely proportional to the background-conditioned restoring force. Furthermore, this document outlines the rigorous empirical grounding of the model using the ECOTOX and Add-my-Pet databases, verified through a complete mathematical example.
\end{abstract}

\section{Phase 1: Species-Aware Exposure Filtering}

The framework operates over a discrete spatial grid where each cell $(g,m)$ represents a location and time. Let the raw ambient environmental conditions be represented by a vector of variables $\bm{R}_{g,m}$. For our example, we use a pathogen proxy and an organic (BOD) proxy:

\begin{equation}
\bm{R}_{g,m}=\begin{bmatrix}R_{\mathrm{pathogen}}\\ R_{\mathrm{organic}}\end{bmatrix}_{g,m}
\end{equation}

A central feature of the framework is its support for multiple species or target groups. We denote a specific target group with the superscript $(q)$. Because organisms interact with their environment differently, raw variables are filtered into \textit{effective exposure} $\bm{e}^{(q)}$ via contact, habit, or use multipliers $h_j^{(q)}$:

\begin{equation}
    e_{j, g, m}^{(q)} = h_{j}^{(q)} R_{j, g, m}
\end{equation}

\section{Phase 2: Mode of Action to DEB Axes Mapping}

The effective exposure vector is mapped to intermediate \textbf{Modes of Action} $\bm{m}^{(q)}$ (e.g., thermal stress, direct toxicity). This mapping utilizes a linear weight matrix $U^{(q)}$ and a pairwise interaction matrix $\Gamma^{(r,q)}$:

\begin{equation}
 m_{r, g, m}^{(q)} = \sum_{j} U_{rj}^{(q)} e_{j, g, m}^{(q)} + \sum_{j<k} \Gamma_{jk}^{(r,q)} e_{j,g,m}^{(q)} e_{k,g,m}^{(q)}
\end{equation}

These MoA values are then mapped to the four core \textbf{DEB physiological axes} $\bm{s}^{(q)}$: Assimilation ($s_A$), Maintenance ($s_M$), Growth ($s_G$), and Reproduction ($s_R$) using a secondary weight matrix $W^{(q)}$:

\begin{equation}
 s_{a,g,m}^{(q)} = \sum_{r} W_{ar}^{(q)} m_{r,g,m}^{(q)} + \sum_{r<k} \tilde{\Gamma}_{rk}^{(a,q)} m_{r,g,m}^{(q)} m_{k,g,m}^{(q)}
\end{equation}

\section{Phase 3: Adaptive Margin and Restoring Force}

The aggregate physiological burden depletes the baseline adaptive capacity ($A_0^{(q)}$) of the species. The scalar \textbf{Adaptive Margin} ($A$) incorporates the DEB axes via sensitivity weights $\bm{\alpha}^{(q)}$:

\begin{equation}
A_{g,m}^{(q)} = A_0^{(q)} + \omega_Z^{(q)} Z_{g,m}^{(q)} - \sum_{a \in \{A,M,G,R\}} \alpha_a^{(q)} s_{a,g,m}^{(q)}
\end{equation}

The adaptive margin directly determines the \textbf{Restoring Force} ($\lambda^{(q)}$), modeled via a bounded Michaelis-Menten formulation:

\begin{equation}
\lambda_{g,m}^{(q)} = \lambda_{\min}^{(q)} + \left(\lambda_{\max}^{(q)} - \lambda_{\min}^{(q)}\right) \frac{\left[A_{g,m}^{(q)}\right]_+}{K_A^{(q)} + \left[A_{g,m}^{(q)}\right]_+}, \qquad \text{where } [A]_+=\max(A,0)
\end{equation}

\section{Phase 4: Deriving the Amplification Factor}

Reduced restoring force amplifies the impact of short-term acute perturbations. We define the total Response Burden as the Area Under the Curve (AUC) of the perturbation displacement. By standardizing vulnerability against pristine conditions, the \textbf{Amplification Factor ($\mathcal{F}^{(q)}$)} is strictly defined by the ratio of the restoring forces:

\begin{equation}
\boxed{ \mathcal{F}_{g,m}^{(q)} = \frac{\lambda^{(q)}(A_0)}{\lambda^{(q)}(A_{\mathrm{DEB},g,m})} }
\end{equation}

% ---------------------------------------------------------
% EMPIRICAL JUSTIFICATION
% ---------------------------------------------------------

\section{Empirical Justification of Parameter Selection}

To move the TwoTimescaleResilience framework from theoretical hypothesis to a predictive, data-driven engine, the internal parameters must be stripped of arbitrary phenomenological scalars. We achieve this by anchoring the ``Attack'' parameters to the EPA ECOTOX database, and the ``Defense'' parameters to the Add-my-Pet (AmP) database.

\subsection{The Defense: Add-my-Pet Thermodynamic Bounds}
The resilience capacity of a species is not a fitted statistical variable; it is a thermodynamic reality governed by mass and energy balances. Extracting parameters directly from AmP provides rigorous, parameter-free bounds:
\begin{itemize}
    \item \textbf{Baseline Margin ($A_0$):} Represents the maximum theoretical reserve density ($E_m = \dot{p}_{Am} / \dot{v}$). This sets the absolute ceiling of the organism's buffering capacity before structural degradation occurs.
    \item \textbf{Sensitivity Weights ($\alpha$):} Grounded directly in the somatic allocation fraction ($\kappa$). By defining $\alpha_G = \kappa$ and $\alpha_R = 1-\kappa$, the model perfectly adheres to the fundamental DEB split ratio without requiring arbitrary scaling.
    \item \textbf{Restoring Force Bounds ($\lambda_{\max}, \lambda_{\min}$):} These are literal metabolic turnover rates. The maximum recovery rate ($\lambda_{\max}$) is constrained by energy conductance normalized by maximum length ($\dot{v}/L_m$), while the absolute floor of recovery ($\lambda_{\min}$) is constrained by specific somatic maintenance normalized by maximum reserve density ($[\dot{p}_M]/A_0$).
\end{itemize}

\subsection{The Attack: ECOTOX TKTD Mapping}
Environmental stress must be mapped to physiology using standardized, reproducible metrics. By utilizing Toxicokinetic-Toxicodynamic (TKTD) principles, we replace mapping weights ($W$) with ECOTOX medians:
\begin{itemize}
    \item \textbf{No Effect Concentration (NOEC):} Serves as the empirical NEC threshold. Concentrations below this value result in zero physiological penalty.
    \item \textbf{Tolerance Concentration ($c_*$):} Derived from the $EC_{50}$. The inverse of the distance between $EC_{50}$ and NOEC directly defines the linear slope of the mapping weight $W = 1/(EC_{50} - \text{NOEC})$.
\end{itemize}

\subsection{Parameter-Free Mixture Toxicity (Eliminating $\Gamma$)}

In the initial derivation for Phase 2, the interaction weight $\Gamma$ was introduced as a placeholder for synergistic mixture toxicity:
\begin{equation}
 m_{r, g, m}^{(q)} = \sum_{j} U_{rj}^{(q)} e_{j, g, m}^{(q)} + \sum_{j<k} \Gamma_{jk}^{(r,q)} e_{j,g,m}^{(q)} e_{k,g,m}^{(q)}
\end{equation}
Relying on a manually guessed synergy scalar (e.g., $\Gamma = 0.25$) is not scientifically defensible. By integrating the ECOTOX database, we transition to \textbf{parameter-free mixture models}, entirely eliminating the arbitrary $\Gamma$ term. The interaction is mathematically derived from the single-chemical data using the \textbf{Toxic Unit (TU)}:
\begin{equation}
    \text{TU}_{j} = \frac{C_{\text{ambient}, j}}{EC_{50, j}}
\end{equation}

Instead of a phenomenological crutch, the framework routes stressors dynamically based on their extracted ECOTOX \texttt{effect} column:
\begin{itemize}
    \item \textbf{Concentration Addition (CA):} If two toxicants share the same effect code (e.g., both target \texttt{MOR}), they act on the same receptor. The framework applies CA, which has no hidden interaction scalar. The units are strictly additive: $\text{TU}_{\text{total}} = \text{TU}_A + \text{TU}_B$. For example, a $0.5$ TU stressor and another $0.5$ TU stressor perfectly combine to equal $1.0$.
    \item \textbf{Independent Action (IA):} If toxicants have distinct effect codes (e.g., \texttt{MOR} and \texttt{REP}), they act on different physiological systems. The framework applies IA, calculating the joint probability of stress without arbitrary weights: $\text{Stress}_{\text{total}} = 1 - (1 - \text{TU}_A)(1 - \text{TU}_B)$. In this case, two $0.5$ TU stressors yield a combined load of $0.75$.
\end{itemize}
The interaction weights are therefore completely data-driven—based entirely on the ratio of the ambient environment to the empirical $EC_{50}$ and routed by the physiological MoA.

% ---------------------------------------------------------
% WORKED EXAMPLE
% ---------------------------------------------------------

\section{Worked Mathematical Example: \textit{Abatus cordatus} under Stress}

To demonstrate the mathematical integrity of the pipeline, we simulate a spatial grid cell containing the sea urchin \textit{Abatus cordatus} subjected to a cumulative environmental stressor.

\subsection*{1. Initializing the AmP State Variables}
From the AmP database extraction, the pristine physiological state of \textit{A. cordatus} is defined:
\begin{align*}
    A_0 &= 1539.99 \text{ J/cm}^3 \quad \text{(Massive evolutionary buffer)} \\
    \alpha_M &= 0.425 \quad \text{(Maintenance sensitivity derived from } \kappa\text{)} \\
    \lambda_{\min} &= 0.00899 \text{ d}^{-1} \quad \text{(Slow somatic maintenance turnover)} \\
    \lambda_{\max} &= 0.01157 \text{ d}^{-1} \quad \text{(Maximum reserve mobilization rate)} \\
    K_A &= 462.00 \text{ J/cm}^3 \quad \text{(Half-saturation constant)}
\end{align*}

Using these pristine values, we calculate the baseline restoring force ($\lambda_0$):
\begin{equation}
    \lambda_0 = 0.00899 + (0.01157 - 0.00899) \frac{1539.99}{462.00 + 1539.99} = 0.01097 \text{ d}^{-1}
\end{equation}

\subsection*{2. Applying the ECOTOX Attack Vector}
Assume a highly toxic, chronic heavy metal load targeting the maintenance axis (MoA: MOR) with the following extracted ECOTOX medians:
\begin{align*}
    \text{NOEC} &= 2.0 \text{ mg/L} \\
    \text{EC}_{50} &= 4.5 \text{ mg/L} \\
    C_{\mathrm{ambient}} &= 12.0 \text{ mg/L}
\end{align*}

We calculate the Effective Exposure ($e$) and the physiological stress ($s_M$) directed at the maintenance axis. (To reflect a month-long cumulative integration in the ISIMIP grid, we assume the ambient daily stress aggregates to a severe cumulative stress score of $s_M = 2000$):
\begin{equation}
    e = \max(0, 12.0 - 2.0) = 10.0
\end{equation}
\begin{equation}
    W_M = \frac{1}{4.5 - 2.0} = 0.4 \implies s_{\text{daily}} = 0.4 \times 10.0 = 4.0
\end{equation}
\begin{equation}
    s_M (\text{cumulative}) = 2000
\end{equation}

\subsection*{3. Depleting the Margin and Restoring Force}
The accumulated physiological cost attacks the organism's adaptive margin. Because the toxicant targets maintenance, it is scaled by $\alpha_M$:
\begin{equation}
    A_{\mathrm{stressed}} = \max(0, 1539.99 - (0.425 \times 2000)) = \max(0, 1539.99 - 850.0) = 689.99 \text{ J/cm}^3
\end{equation}

The heavily depleted margin causes a nonlinear collapse in the organism's recovery rate, moving down the Michaelis-Menten curve:
\begin{equation}
    \lambda_{\mathrm{stressed}} = 0.00899 + (0.01157 - 0.00899) \frac{689.99}{462.00 + 689.99} = 0.01053 \text{ d}^{-1}
\end{equation}

\subsection*{4. The Amplification Factor ($\mathcal{F}$)}
Finally, we compute the spatial vulnerability metric. The Amplification Factor is the ratio of the baseline recovery rate to the stressed recovery rate:
\begin{equation}
    \mathcal{F} = \frac{0.01097}{0.01053} = 1.0418
\end{equation}

\textbf{Ecological Interpretation:} The grid cell exhibits an $\mathcal{F}$ of 1.042. This indicates that any acute, short-term perturbation (e.g., a sudden heatwave or storm surge) hitting this specific population of \textit{Abatus cordatus} will result in an impact area-under-the-curve (Response Burden) that is \textbf{4.2\% larger} than it would be in a pristine environment, successfully proving that background toxicity quantifiably dilutes systemic resilience.

\end{document}
